% Junior Level 3, Week 4 -- answer key
\answersweek{4}{How Do Computers Learn?}{Answers for workbook chapter 4 and play-at-home card 4.}

\section{Lesson questions}
\begin{answer}[Exit question]
The robot might guess \enquote{not a cat}, because every cat it learned from was orange, so it may think \enquote{cat} means \enquote{orange}. Omar can fix it by showing it cats of many colors (black, white, gray, striped). Answers like \enquote{show it more different cats} show understanding.
\end{answer}

\begin{answer}[Reflection]
The Robot learned from examples (YES and NO cards), not from being told the rule. Sam's Robot made mistakes because the examples were one-sided: every bird it saw could fly. Any drawing of a made-up creature is right; ask which features it has.
\end{answer}

\begin{answer}[Safety question]
Do not answer. Close the chatbot (or stop typing) and tell a grown-up. A name and a home address are personal information.
\end{answer}

\section{Workbook exercises}
\begin{answer}[Exercise 4.1: Teach the Robot]
Any tally is right if it matches the guesses made. Features that fit each rule card: HAS FOUR LEGS: \textbf{legs}. IS A BIRD: \textbf{wings} (feathers and a beak also count, though they are not in the list). CAN FLY: \textbf{wings}. If the Robot circles fur for \enquote{has four legs}, ask about the frog and the hippo: they have four legs but little or no fur.

\textbf{What to look for:} \tickbox{} the Trainer gave no hints \quad \tickbox{} the Robot named a feature, not only a card \quad \tickbox{} the fifth tally mark goes across.
\end{answer}

\begin{answer}[Exercise 4.2: Robot Report]
The two bars match the tally from 4.1: one square for each right guess and one for each wrong guess, both bars starting at the bottom line. A Robot often gets a few wrong at first and more right after more examples.
\end{answer}

\begin{answer}[Exercise 4.3: The One-Sided Robot]
The Robot guesses \textbf{not a bird}. Every YES example it saw was a bird that flies, so it learned \enquote{bird means flying}, and the kiwi does not fly. (It would also call a flying bug a bird.) Sam can fix it with more and different examples: put the kiwi (or a penguin, an ostrich, a chicken) on YES, and the bug on NO.
\end{answer}

\section{Play-at-home card}
\begin{answer}[Teach your grown-up]
Any rule is fine. Look for: the child says only YES or NO, does not give the rule away, and counts the right guesses. With only 2 examples, the Robot usually makes more mistakes: fewer examples, worse guesses.
\end{answer}

\begin{answer}[Safe or not safe?]
\enquote{What is your favorite color?}: fine to answer in a game with a grown-up. \enquote{What school do you go to?} and \enquote{What is your address?}: stop and tell. These are personal information, and we never tell them to a chatbot.
\end{answer}
